\documentclass[a4paper, twocolumn]{article}

\usepackage[utf8]{inputenc}
\usepackage[T1]{fontenc}
\usepackage[slovene]{babel}
\usepackage{amsmath}
\usepackage{amssymb}
\usepackage{graphicx}
\usepackage{geometry}
\usepackage{hyperref}
\usepackage{microtype}
\usepackage{float}
\usepackage{caption}
\captionsetup{skip=2pt}
\raggedbottom

\geometry{margin=2cm, top=2.5cm}

\title{\textbf{Ničle Airyjeve funkcije}}
\author{Sinja Kočica\\
\small Univerza v Ljubljani, Fakulteta za računalništvo in informatiko}
\date{}

\begin{document}

\maketitle

\section{Opis naloge}

Airyjeva funkcija $\mathrm{Ai}(x)$ je rešitev diferencialne enačbe drugega reda:
\[
\mathrm{Ai}''(x) = x\,\mathrm{Ai}(x),
\]
z začetnima pogojema pri $x = 0$:
\[
\mathrm{Ai}(0) = \frac{1}{3^{2/3}\,\Gamma\!\left(\tfrac{2}{3}\right)}, \quad
\mathrm{Ai}'(0) = -\frac{1}{3^{1/3}\,\Gamma\!\left(\tfrac{1}{3}\right)}.
\]
Za negativne vrednosti $x$ funkcija oscilira in večkrat seka os. Naloga zahteva iskanje čim več ničel funkcije $\mathrm{Ai}(x)$ na 10 decimalk natančno, brez uporabe vgrajenih reševalnikov diferencialnih enačb.

\section{Opis rešitve}

Rešitev poteka v dveh korakih. Najprej z Magnusovo metodo numerično izračunamo vrednosti Airyjeve funkcije na gosti mreži točk vzdolž negativne osi, kjer ima funkcija ničle. Na mreži poiščemo intervale, kjer funkcija zamenja predznak, nato pa z metodami za iskanje ničel natančno poiščemo mesta, kjer funkcija preide skozi absciso.

\subsection{Reševanje ODE z Magnusovo metodo}

Enačbo drugega reda prevedemo v sistem prvega reda z uvedbo vektorja stanja $\mathbf{y}(x) = [\mathrm{Ai}(x),\, \mathrm{Ai}'(x)]^\top$:
\[
\mathbf{y}'(x) = A(x)\,\mathbf{y}(x), \quad
A(x) = \begin{pmatrix} 0 & 1 \\ x & 0 \end{pmatrix}.
\]
Za numerično integracijo uporabimo \emph{Magnusovo metodo reda 4}. Na vsakem koraku dolžine $h$ izračunamo vrednosti matrike $A$ v dveh Gaussovih točkah $c_1$ in $c_2$, ki sta izbrani tako da je aproksimacija točna za polinome do stopnje 3:
\[
c_1 = x_k + \left(\tfrac{1}{2} - \tfrac{\sqrt{3}}{6}\right)h, \quad
c_2 = x_k + \left(\tfrac{1}{2} + \tfrac{\sqrt{3}}{6}\right)h.
\]
Iz vrednosti $A_1 = A(c_1)$ in $A_2 = A(c_2)$ sestavimo Magnus eksponent $\sigma_{k+1}$, ki vsebuje dva člena. Prvi člen $\frac{h}{2}(A_1 + A_2)$ je uteženo povprečje matrike $A$ v obeh točkah in bi zadostoval, če bi bila matrika vzdolž koraka konstantna. Ker pa se $A(x)$ spreminja in matrično množenje ni komutativno, samo povprečje ni dovolj. Komutator $[A_1, A_2] = A_1 A_2 - A_2 A_1$ izmeri koliko se matriki $A_1$ in $A_2$ razlikujeta, drugi člen $\frac{\sqrt{3}}{12}h^2[A_1, A_2]$ pa doda potreben popravek:
\[
\sigma_{k+1} = \frac{h}{2}(A_1 + A_2) - \frac{\sqrt{3}}{12}h^2[A_1, A_2].
\]
Nov vektor stanja dobimo kot:
\[
\mathbf{y}_{k+1} = e^{\sigma_{k+1}}\,\mathbf{y}_k.
\]
Postopek ponavljamo vzdolž celotne mreže in s tem dobimo vrednosti $\mathrm{Ai}(x)$ in $\mathrm{Ai}'(x)$ v vsaki točki.

\subsection{Iskanje ničel}

Ko imamo vrednosti $\mathrm{Ai}(x)$ na gosti mreži, poiščemo intervale kjer funkcija zamenja predznak, tam se nahaja ničla. Za vsak tak interval nato uporabimo dve metodi:

\textbf{Bisekcija} na vsakem koraku razpolovi interval $[a, b]$ in obdrži tisto polovico v kateri je ničla:
\[
c = \frac{a + b}{2}, \quad \text{nato } [a,c] \text{ ali } [c,b].
\]
Konvergira zanesljivo, a počasi. Napaka se z vsako iteracijo razpolovi.

\textbf{Newtonova metoda} na vsakem koraku izračuna tangento na krivuljo in njeno presečišče z osjo vzame kot naslednji približek:
\[
x_{k+1} = x_k - \frac{\mathrm{Ai}(x_k)}{\mathrm{Ai}'(x_k)}.
\]
Vrednost odvoda $\mathrm{Ai}'(x)$ dobimo neposredno iz Magnusove metode, saj ta hkrati izračuna oba elementa vektorja $\mathbf{y}$. Newtonova metoda konvergira kvadratično. Vsak korak podvoji število pravilnih decimalk.

\section{Primeri uporabe}

V nadaljevanju prikazujemo delovanje implementacije na konkretnem primeru. To zajema izris Airyjeve funkcije, iskanje ničel in primerjavo metod.

\subsection{Airyjeva funkcija}

Slika~\ref{fig:airy} prikazuje Airyjevo funkcijo in njen odvod na intervalu $[-12, 1.5]$, izračunana z Magnusovo metodo s korakom $h = 0.005$. Za negativne $x$ funkcija oscilira, za pozitivne pa eksponentno pada proti nič. Vse ničle se nahajajo na negativni polosi.

\begin{figure}[H]
    \centering
    \includegraphics[width=\linewidth]{demo/airy_funkcija.png}
    \caption{Airyjeva funkcija $\mathrm{Ai}(x)$ in njen odvod.}
    \label{fig:airy}
\end{figure}

\subsection{Ničle Airyjeve funkcije}

Z Magnusovo metodo in bisekcijo smo poiskali prvih 10 ničel Airyjeve funkcije. Tabela~\ref{tab:nicle} prikazuje izračunane vrednosti skupaj z napakami glede na referenčne vrednosti. Vse napake so manjše od $4 \times 10^{-11}$, kar ustreza zahtevani natančnosti $10^{-10}$.

\begin{table}[H]
\centering
\caption{Prvih 10 ničel $\mathrm{Ai}(x)$ in napake glede na DLMF.}
\label{tab:nicle}
\small
\begin{tabular}{rll}
\hline
$k$ & Ničla $a_k$ & Napaka \\
\hline
1  & $-2.3381074104$ & $1.4 \times 10^{-11}$ \\
2  & $-4.0879494441$ & $9.7 \times 10^{-13}$ \\
3  & $-5.5205598281$ & $4.8 \times 10^{-13}$ \\
4  & $-6.7867080901$ & $9.9 \times 10^{-12}$ \\
5  & $-7.9441335871$ & $4.0 \times 10^{-11}$ \\
6  & $-9.0226508533$ & $1.4 \times 10^{-12}$ \\
7  & $-10.0401743416$ & $1.1 \times 10^{-11}$ \\
8  & $-11.0085243037$ & $1.4 \times 10^{-12}$ \\
9  & $-11.9360155632$ & $1.1 \times 10^{-12}$ \\
10 & $-12.8287767529$ & $9.1 \times 10^{-12}$ \\
\hline
\end{tabular}
\end{table}

\subsection{Primerjava metod za iskanje ničel}

Slika~\ref{fig:konvergenca} prikazuje konvergenco bisekcije in Newtonove metode pri iskanju prve ničle. Bisekcija potrebuje okoli 35 iteracij, Newton pa le 3. V povprečju je Newton $8.8\times$ hitrejši.

\begin{figure}[H]
    \centering
    \includegraphics[width=\linewidth]{demo/konvergenca.png}
    \caption{Konvergenca bisekcije in Newtonove metode.}
    \label{fig:konvergenca}
\end{figure}

\begin{figure}[H]
    \centering
    \includegraphics[width=\linewidth]{demo/primerjava_iteracij.png}
    \caption{Število iteracij za prvih 10 ničel.}
    \label{fig:iteracije}
\end{figure}

\subsection{Vpliv koraka na natančnost}

Manjši korak $h$ v Magnusovi metodi ne pomeni vedno večje natančnosti. Pri zelo majhnem koraku se nakopičijo zaokrožitvene napake, ki začnejo prevladovati nad napako metode same. Slika~\ref{fig:napake_h} prikazuje primerjavo napak za $h = 0.0005$ in $h = 0.0001$. Pri nekaterih ničlah je manjši korak celo slabši, kar potrjuje da je $h = 0.0005$ bila dobra izbira.

\begin{figure}[H]
    \centering
    \includegraphics[width=\linewidth]{demo/napake_h.png}
    \caption{Napake ničel za različni vrednosti koraka $h$.}
    \label{fig:napake_h}
\end{figure}

\end{document}
